\scanpage{20}
% archive foliation: 19
Plus généralement \deleted{on} posons
\[
\mathcal H^q(\Phi)(S)=H^q(-/S,\Phi)=\left\{\begin{gathered}\text{valeur en }\Phi\text{ du}\\\text{foncteur dérivé qième}\\\text{de }\Phi\longrightarrow\Phi(S)\end{gathered}\right.
\]
\margin{\unsure{identités} à développer $\longrightarrow$}
\inserted{définition \unsure{valeur en} $\Phi|S$ \unsure{du foncteur dérivé} de $\Gamma^S$ sur la catégorie des \deleted{\unsure{hyperfaisceaux}} \unsure{hyperfaisceaux} au dessus de $S$ ; \unsure{il en résulte immédiatement}}

\figure{\begin{tikzpicture}\draw (0,0) to[out=-100,in=100] (0.12,-0.55) to[out=-80,in=120] (0.5,-2.7);\end{tikzpicture}}
On sait que

(i) $\mathcal H^q(\Phi)$ est un hyper-préfaisceau qui \deleted{\unsure{transforme}} transforme sommes \inserted{finies} en produits

(ii) $\mathcal H^q(\Phi_A)(G/F)=H^q(F,A)$

Sous ces conditions, on a la suite spectrale suivante, si $T\longrightarrow S$ est \uline{surjectif} (dans $C_G$)
\[
\text{\textcircled{B}}\qquad\boxed{H^*(-/S,\Phi)\Longleftarrow H^p(T/S,\mathcal H^q(\Phi))}
\]
\margin{\fbox{\parbox{3cm}{suite spectrale de Leray - Čech}}}

Si $F'\subset F\subset G$, on peut prendre $T=G/F'\qquad S=G/F$ et on obtient, si $A$ est un $G$-module
\[
H^*(F,\deleted{\unsure{\Phi}}A)\Longleftarrow H^p(F^{\unsure{\natural|n}}_{\mathrm{mod}\,F'};\mathcal H^q(\deleted{\unsure{\Phi}}A))
\]
La suite spectrale est la $=$ que si on oublie sur $A$ la structure de $G$-module, et on retient que celle de $F$-module. La suite spectrale \textcircled{B} n'est donc pas essentiellement plus générale que \textcircled{A}.

\margin{$\Vert$ Tout ceci est à généraliser aux hyperfaisceaux quelconques \quad -- -- -- -- --}
